\chapter{}

Un altro anno; e s’era alla fine di marzo, il giorno di San Giuseppe.

Da Poggio a’ Meli s’udivano gli scampanii, che si rimescolavano alla rinfusa nel cielo come un suono che crescesse sempre, quasi immobile, con una romba greve. E a Pietro era venuta un’allegria insolita, un’allegria simile ad un benessere troppo forte, che lo faceva più nervoso.

Vorrei parlare di questi indefinibili turbamenti del marzo, a cui è unita quasi sempre una sottile voluttà, un desiderio di qualche bellezza.

Questi soli ambigui, questi cinguettii ancora nascosti e che si dimenticano presto, queste nuvole biancheggianti che sembrano venute prima del tempo! E le foglie secche, che sono ancora sopra i grani germogliati, mescolando il pallore della morte con il pallore della vita! Queste foglie di tutte le specie, che si trovano ancora sopra l’erbe per rinnovarsi; le piante potate, e i loro rami e i loro tralci, sparsi a terra, che saranno portati via per sempre! E questi rami secchi tagliati dai frutti, che esitano ancora a fiorire su le rame nuove! La terra un poco umida, che s’attacca alla punta delle vanghe, e i contadini sono costretti a pulirle con i pollici; e le zolle che rimangono agli zoccoli di legno! E quest’amore quasi matrimoniale e sconosciuto a noi di tutti gli esseri che s’aiutano; e anche i loro odii! E il vischio che nasce su i rami dei testucchi, tagliato con un colpo di pennato! Ma farà subito il ributto. E le gemme dei castagni!

Domenico andò nel campo, seguito dai suoi assalariati, per combinare le faccende dell’indomani.

Pietro era grasso, ma pallido e con un’aria di gracilità: entrava in quindici anni. Credeva che fosse ridicola e disadatta alla sua età la giubba con il bavero alla marinaia, tagliata per economia da una veste vecchia.

Entrò svelto in casa di Giacco; il quale, come il solito, gli mise una mano su la spalla:

— Come cresce a fretta! Scommetto che mi ha portato da fumare.

Pietro gli prese i baffi e glieli tirò di qua e di là; Giacco per non sentir male era pronto a girare il collo.

Il ragazzo rise, guardando Masa, che disse:

— Più forte.

— No, no; ora basta.

E lo allontanò da sè a poco a poco, ma risolutamente. Poi chiese:

— Dunque, nè meno una cicca?

Rebecca, spazzando la trattoria, metteva in serbo le cicche trovate, e lo incaricava di portargliele.

Masa intervenne un’altra volta:

— Non fuma mica il padroncino!

E ne rise insieme con lui come di una burla. Dopo avere riso, storceva le labbra e se le mordeva. Il vecchio cavò dal taschino una pipa sbocconcellata, con una cannuccia corta quanto il palmo della mano.

— Grazie a Dio, ci ho sempre quello che la sua mamma mi dette la settimana passata. Guardi se non è vero!

Battè la pipa in proda alla tavola: schizzò fuori una specie di polvere incenerita. Egli la radunò insieme, la mescolò e la rimise dentro. Poi prese, dal focolare, un fuscello acceso. A stento, gli uscì di bocca un poco di fumo, azzurro chiaro. Ed egli, guardandolo, disse:

— Oh, c’è poco trinciato, oggi!

Indi con il pollice che aveva l’unghia mozzata da un taglio fattosi da giovine, pigiò dentro il pezzetto di brace rimasta nella pipa.

Pietro vide un’altra volta quel fumo, e, dentro di sè, come una cosa reale, che gli dette un malessere, la mamma che andava a un cassetto, in casa, e voleva prendere qualche cosa. Ma tutti s’erano allontanati da lei! E mentr’ella si ostinava, il cassetto spariva nel muro. Allora gli parve di sentire sul volto le sue mani, come un grande bacio, come se le mani lo baciassero.

Masa, meravigliata della sua espressione sbigottita gli chiese:

— A che pensa?

Il vecchio si avvicinò all’uscio, e disse:

— Bisogna che vada a governare le vacche. Dammi la fune.

Ma Masa, preoccupata di vedere il padroncino così, rispose di malumore:

— Dove l’hai messa?

E Giacco:

— Cercamela.

— Non sai mai quello che fai. Poi ti ci vuole la moglie intorno per darti quello che ti manca.

— Quanto chiacchieri! Se tu avessi trovato la fune, senza rispondermi niente? Non avresti fatto meglio?

— Io chiacchiero quanto mi pare; quanto te.

Poi chiese a Pietro, per distrarlo; credendo che soffrisse di qualche rimbrotto:

— Ha visto Ghìsola oggi?

Rispose egli sbadatamente:

— Non è qui?

— È voluta andare alla messa a Siena.

Disse Giacco, con l’aria di chi ripiglia un battibecco. Ma Masa la difese:

— Ha fatto bene. Qui a Poggio a’ Meli non si vede mai nessuno.

E a Pietro soggiunse:

— Credevo che l’avesse incontrata!

I due vecchi divennero pensierosi, guardandosi con occhiate che Pietro non comprendeva. Masa esclamò sospirando:

— Sarà quel che Dio vuole!

— Di che cosa? — chiese Pietro. — Ditemelo.

Un’acre curiosità lo invase:

— Ma dov’è? Tornerà tra poco?

Si sentì sbigottito; e si vide subito dai suoi occhi azzurri, sempre così buoni che tutti lo sapevano: le palpebre gli sembrarono come acqua calda.

Il cavallo attaccato al calesse, legato nel piazzale ad una campanella di ferro, si ripiegava tutto da una parte, riposandosi. Toppa finiva un seccarello terroso; tenendolo fermo con le zampe; per roderlo meglio.

Pietro non era ancora calmo quando scorse Ghìsola.

Era divenuta una giovinetta. I suoi occhi neri sembravano due olive che si riconoscono subito nella rama, perchè sono le più belle; quasi magra, aveva le labbra sottili.

Egli si sentì esaltare: ella camminava adagio smuovendo un poco la testa, i cui capelli nerissimi, lisciati con l’olio, erano pettinati in modo diverso da tutte le altre volte.

Cercò di smettere il suo sorriso, abbassando il volto; ma rallentò il passo, come se fosse indecisa a voler dissimulare qualche segreto. Egli ne ebbe un dispiacere vivo, e le mosse incontro, come quando più ragazzi, per farle un dispetto oppure per raccontarle qualche cosa, con la voglia d’offenderla.

Come s’era imbellita da che non l’aveva più veduta!

Notò, con gelosia, un nastro rosso tra i suoi capelli, le scarpe lustre di sugna e un vestito bigio quasi nuovo; e fece un sospiro.

Ma ella, così risentita che non gli parve nè meno possibile, gli gridò:

— Vada via, c’è suo padre. Non mi s’avvicini.

Egli, invece, continuò ad andarle incontro; ma ella fece una giravolta, rasentandolo senza farsi toccare. Pietro non le disse più niente, non la guardò nè meno: era già offeso e mortificato. Perchè si comportava così? Sarebbe andato a trovarla anche in casa, dov'ella entrò soffermandosi prima con un piede su lo scalone! Si struggeva; era assillato da una cosa che non comprendeva; aveva voglia d'imporlesi.

Ma, a poco a poco, si sentì rappacifichito e lieto un'altra volta; come se non le dovesse rimproverare nulla; mentre un sentimento delizioso gli si affermava sempre di più.

Ghìsola riuscì presto di casa: s'era tolto il nastro, aveva cambiato le scarpe, mettendosi un grembiule rosso sbiadito. Alzò gli occhi verso Pietro, seria e muta; ed entrò in capanna dimenandosi tutta. Pose dentro una cesta il fieno già falciato dal nonno; poi smise, per levarsi una sverza da un dito. Egli si sentì uguale a quella mano. E il silenzio di lei, inspiegabile, lo imbarazzò; e non sarebbe stato capace a parlare per primo. Perciò le dette una spinta, ma lieve; ed ella, fingendo d'esser stata per cadere, lo guardò accigliata.

Egli disse:

— Quest'altra volta ti butto in terra da vero!

— Ci si provi!

Quand'ella voleva, la sua voce diveniva dura e aspra, strillava come una gallina. Allora egli la guardò con dispetto, sentendo che doveva obbedire.

Per solito, mentre parla, non si sente il suono della voce di chi si ama; o, almeno, non si potrebbe descrivere.

Ella aggiunse:

— Vada via.

Egli provava lo stesso effetto di quando siamo sotto l'acqua e non si possono tenere gli occhi aperti; ma rispose:

— Ghìsola, tu mi dicesti un mese fa che mi volevi bene. Non te ne ricordi? Io me ne ricordo, e ti voglio bene.

E rise, terminando con un balbettìo. Ghìsola lo guardò come se ci si divertisse; e, in fatti, le piacque quel ripiego d'inventare una cosa per dirne una vera.

Ella rispose:

— Lo so, lo so.

Egli, invece di poter seguitare, notò come la tasca del suo grembiule era graziosa. E di lì, d'un tratto, le tolse il piccolo fazzoletto orlato, alla meglio, di stame celeste.

— Me lo renda.

Egli, temendo di aver fatto una sciocchezza, glielo rese.

— Ti sei bucata codesto dito?

Riuscendo a parlare, non gli parve poco.

— Che cosa le importa? Tanto lei non lavora. Non fa mai niente.

Gli rispose con superbia burlesca e sfacciata; ma egli la prese sul serio e disse:

— Ghìsola, se vuoi, ti aiuto.

Ella finse di canzonarlo come se non fosse stato capace; e lo allontanò dicendogli che non voleva aiutarla, ma toccarla.

Domenico sopraggiunse dal campo.

Pietro raccolse in fretta un olivastro, ch'era lì in terra; e cominciò a frustare l'uscio della capanna come per uccidere le formiche, che lo attraversavano in fila.

Ghisola si chinò a prendere a manciate il fieno, con movimenti bruschi e rapidi; e, voltasi dalla parte del mucchio, finì d'empire la cesta. Poi l'alzò per mettersela in spalla, ma non fu capace da sé: gli ossi dei bracci pareva che le volessero sfondare i gomiti.

Allora Pietro l'aiutò prima che il padre potesse vedere. Ghìsola, assecondando il movimento di lui, guardava verso Domenico con i suoi occhi acuti e neri, quasi che le palpebre tagliassero come le costole di certi fili d'erba. Ma Pietro arrossì e tremò perchè ella, innanzi di muovere il passo, gli prese una mano. Rimase sbalordito, con una tale dolcezza, che divenne quasi incosciente; pensando: «Così dev'essere!».

Domenico, toccati i finimenti del cavallo se erano ancora affibbiati bene, gli gridò:

— Scioglilo e voltalo tu. Ripiega la coperta e mettila sul sedile.

La bestia non voleva voltare; e lo sterzo delle stanghe restava a traverso. Anche lo sguardo di Toppa, sempre irato, molestava e impacciava Pietro.

— Tiralo a te!

Non aveva più forza, non riesciva ad afferrare bene la briglia; e le dita gli entravano nel morso bagnato di bava verdognola e cattiva. Nondimeno fece di tutto, anche perchè sapeva che Ghìsola, tornata dalla stalla, doveva essere lì. Tremava sempre di più. E le zampe del cavallo lo rasentarono, poi lo pestarono. Allora Domenico prese in mano la frusta, andò verso Pietro e gliel'alzò sul naso.

— Lo so io che hai. Ma ti fo doventare buono a qualche cosa io.

Ghìsola si avvicinò al calesse e lo aiutò; dopo aver sdrusciato, allo spigolo del pozzo, uno zoccolo a cui s'era attaccato il concio della stalla.

Domenico, sempre con la frusta in mano, andò a parlare a Giacco che ascoltava con le braccia penzoloni e i pollici ripiegati tra le dita, le cui vene sollevavano la pelle, come lombrici lunghi e fermi sotto la moticcia.

Pietro non aveva il coraggio di guardare in volto Ghìsola, i cui occhi adesso lo seguivano sempre. Le gambe gli si piegavano, con una snervatezza nuova; che aumentava la sua confusione simile a una malattia. Ghìsola lo aiutò ancora; e, nel prendere la coperta rossa che era stata stesa sul cavallo, le sue dita lo toccarono; nel metterla sul sedile, le loro nocche batterono insieme; ed ambedue sentirono male, ma avrebbero avuto voglia di ridere.

Domenico salì sul calesse, sbirciò Pietro e gridò ancora:

— Sbrigati! Che cos'hai nel labbro di sotto? Pulisciti.

Egli, impaurito, rispose:

— Niente.

Poi pensò che ci fosse il segno delle parole dette a Ghìsola. Ma subito dopo gli dispiacque di essere così sciocco; mentre il cuore gli balzava come per escire fuori.

Gli assalariati e Giacco salutarono, togliendosi il cappello. Pietro a pena ebbe tempo di far con l'angolo della bocca un piccolo cenno a Ghìsola; ma ella era così attenta al padrone che aggrottò in fretta le sopracciglia. Allora Pietro guardò la testa del cavallo, che già tirava il calesse fuori del piazzale mettendosi a trotto a pena nella strada.

La luce del sole tramontato dietro la Montagnola, più rossa che rosea, era sopra a Siena. Ma i cipressi sparsi da per tutto, a fila o a cerchio in cima alle colline, gli dettero il rammarico di staccarsi da una cosa immensa.

Domenico, guidando, non parlava mai; rispondendo con il capo a coloro che lo salutavano. Sorrideva in vece a qualche ragazza che conosceva; e, facendo prima rallentare il cavallo, la toccava con la punta della frusta nel mezzo del grembiule. E Pietro, con gli occhi socchiusi, si voltava dalla parte opposta, arrossendo; poi si distraeva guardando le gambe del cavallo; e gli pareva che il loro rumore variasse di tempo a seconda delle arie che gli passavano per la mente. Oppure cercava di non sentire quell'odore particolare, che avevano gli abiti del padre.
